% Supplementary material: audit protocol detail moved from Section 5.1 of the main text.

\section{Audit Protocol Detail}
\label{supp:audit}

This section expands the protocol summary of Section~\ref{subsec:reported_protocol} of the main text.

\subsection{Benchmark Tasks and Metrics}
\label{subsec:reported_benchmarks}
The benchmarks that recur across the reviewed release documents span seven
capability areas: \gls{mmlu} \citep{hendrycks2021mmlu} for broad knowledge;
ARC-Easy and ARC-Challenge \citep{clark2018arc} and HellaSwag
\citep{zellers2019hellaswag} for commonsense and language reasoning; \gls{gsm8k}
\citep{cobbe2021gsm8k} and MATH \citep{hendrycks2021math} for mathematical
reasoning; HumanEval \citep{chen2021humaneval} and MBPP \citep{austin2021mbpp}
for code generation; TriviaQA \citep{joshi2017triviaqa} and SQuAD~2.0
\citep{rajpurkar2018squad2} for question answering; MNLI \citep{williams2018mnli}
for natural language inference; and SST-2 \citep{socher2013sst} for sentiment
classification. Each benchmark carries a standard metric: accuracy for multiple
choice, exact match or F1 for question answering, pass@1 for code, exact match
for \gls{gsm8k}, and symbolic equivalence for MATH. Instruction-following
(MT-Bench-style \citep{zheng2023mtbench}) and summarization are reported too
sparsely and too heterogeneously to compare at all, and the conventions behind
them are described in \suppsec~\ref{subsec:appendix_human_eval}. Section~\ref{subsec:reported_agentic_gap}
takes up separately the question of whether any of these benchmarks bears on
agentic suitability.

Five conditions determine what a benchmark number means, and a source that omits
any of them has not reported a comparable result. \emph{Benchmark version:}
\gls{mmlu} exists in the original and in MMLU-Pro, ARC in Easy and Challenge
splits, GSM8K in original and GSM8K-Platinum revisions, and HumanEval in the
original and HumanEval+ forms, and these are not interchangeable.
\emph{Prompt template:} whether the model's own chat template is applied,
whether options are lettered or restated, and whether a system prompt is present.
\emph{Answer extraction:} whether a multiple-choice answer is scored by
comparing per-option log-likelihoods or by parsing a generated letter, which can
differ by several points on the same model and benchmark, and for
chain-of-thought tasks, which span of the generation is taken as the answer.
\emph{Decoding settings:} greedy versus sampled, temperature, and for pass@$k$
the number of samples. \emph{Harness:} the framework and its version, most often
\texttt{lm-evaluation-harness} \citep{gao2023harness} or HELM
\citep{liang2023helm}, each with its own defaults for all of the above.

The coverage of these five conditions across the reviewed sources is what
Section~\ref{subsec:reported_overall} reports on, and it is poor. Which
benchmarks a vendor reports, in which version, under how many shots, and with
which extraction rule varies enough to prevent a common table.

\subsection{Coverage Ledger}
\label{subsec:audit_ledger}

Table~\ref{tab:audit_ledger} reproduces the coverage ledger summarized in
Section~\ref{subsec:reported_overall}. Its scope is the instruction-tuned
checkpoints of the 14 cohort models from Microsoft, Alibaba, Meta, and Google,
against the eight benchmarks that recur in their headline reporting; the other
four cohort models (TinyLlama, SmolLM2, StableLM~2, Mistral~7B), each from a
different developer, were not transcribed.
Each cell was transcribed from the model's primary release document or official
model card, listed with its table in the released CSV. Only coverage and stated
protocol are recorded, not scores. Of the 112 cells, 69 carry an
instruction-tuned value and 45 of those state a protocol; no benchmark is reported
for all 14 models; and the largest group of models sharing one stated protocol on
any benchmark is five (\gls{mmlu} at 5 shots; \gls{gsm8k} at 8 shots with
chain-of-thought). Of the reconciliation rules written before collection, rule~(a)
(exclude a value that cannot be cross-checked against a second source) would
exclude nearly every vendor cell by construction, so the finding is stated without
it: protocol mismatch alone (rule~(b)) and separate pre-trained and
instruction-tuned tabulation (rule~(c)) are sufficient to prevent a common table.

% Generated by source/scripts/build_audit_ledger.py; do not edit by hand.
\begin{table}[t]
\centering
\caption{Coverage ledger for the instruction-tuned checkpoints of the 14 cohort models from the four largest vendors: the protocol under which each model's primary release document reports each benchmark. Numbers are shot counts; ``c'' marks chain-of-thought; ``SE'' OpenAI simple-evals defaults; ``u'' a value reported without a stated protocol; ``--'' not reported. Sources and locations are listed in \file{audit\_coverage\_ledger.csv}; transcribed on 8~October 2026.}
\label{tab:audit_ledger}
\footnotesize
\begin{tabular}{lcccccccc}
\toprule
\textbf{Model} & \textbf{MMLU} & \textbf{MMLU-Pro} & \textbf{ARC-C} & \textbf{HellaSwag} & \textbf{GSM8K} & \textbf{MATH} & \textbf{HumanEval} & \textbf{MBPP} \\
\midrule
Phi-3-mini & 5 & -- & 10 & 5 & 8c & -- & 0 & 3 \\
Phi-3.5-mini & 5 & 0c & 10 & 5 & 8c & 0c & 0 & 3 \\
Phi-4-mini & 5 & 0c & 10 & 5 & 8c & 0c & -- & -- \\
Phi-4 & SE & -- & -- & -- & -- & SE & SE & -- \\
Qwen2.5-0.5B & -- & u & -- & -- & u & u & u & u \\
Qwen2.5-1.5B & -- & u & -- & -- & u & u & u & u \\
Qwen2.5-3B & -- & u & -- & -- & u & u & u & u \\
Qwen2.5-7B & -- & u & -- & -- & u & u & u & u \\
Llama 3.2-1B & 5 & -- & 0 & 0 & 8c & 0c & -- & -- \\
Llama 3.2-3B & 5 & -- & 0 & 0 & 8c & 0c & -- & -- \\
Gemma 2-2B & u & -- & -- & -- & -- & -- & -- & u \\
Gemma 2-9B & u & -- & -- & -- & -- & -- & -- & u \\
Gemma 3-1B & 0 & 0 & -- & -- & 0c & 0 & 0 & 3 \\
Gemma 3-4B & 0 & 0 & -- & -- & 0c & 0 & 0 & 3 \\
\midrule
Reported (of 14) & 10 & 8 & 5 & 5 & 11 & 11 & 9 & 10 \\
\bottomrule
\end{tabular}
\end{table}


\subsection{Hardware and Quantization Scope}
\label{subsec:reported_hardware}
The review reads reported behavior against six reference configurations spanning the
cloud-to-edge spectrum (Table~\ref{tab:hardware_configs}) and three precisions:
FP16 (reference), INT8, and INT4, with \gls{gguf}~Q4\_K\_M, GPTQ group-128, and \gls{awq}
group-128 treated as distinct INT4 subclasses when the source distinguishes them.
Inference-framework attribution (vLLM \citep{kwon2023vllm}, llama.cpp
\citep{gerganov2023llamacpp}, ONNX Runtime \citep{microsoft2019onnx}, ExecuTorch
\citep{meta2024executorch}, MLC-LLM \citep{mlc2024}) follows the stack used by the
source of each figure; cross-framework comparisons are made only when compared
figures share a framework or are explicitly flagged.

\subsection{Limitations of the Evidence Base}
\label{subsec:reported_limits}
Because the figures are compiled from heterogeneous public sources rather than a
single controlled pipeline, four limitations bound the strength of every
comparison below. \textbf{Prompt-format and harness variation:} different sources
use different templates and answer-extraction heuristics, which can shift reported
accuracy by several points under the same nominal benchmark. \textbf{Self-reporting
bias:} vendor-reported scores are systematically more favorable than third-party
reproductions; where third-party figures exist, comparisons preferentially use
them and flag the divergence. \textbf{Condition-coverage sparsity:} not every
model has reported scores for every task, tier, or precision. Missing values are
never imputed, interpolated, or carried across from a sibling model size or from
a base checkpoint to its instruction-tuned variant; a tuple with no qualifying
source is recorded as absent, and absence is itself reported, since the pattern
of what vendors decline to publish is part of the finding of
Section~\ref{subsec:reported_overall}. No statistic in this paper is computed
over a set from which absent tuples were silently dropped. \textbf{Snapshot staleness:} the ecosystem evolves monthly,
and figures reflect the state of public reporting at the collection window. These
limitations motivate the empirical-verification agenda of
Section~\ref{subsec:challenges_verification}.

\subsection{Quantization Impact}
\label{subsec:reported_quantization}

The reported picture for the sub-10B regime is summarized in
Table~\ref{tab:compression_quality} of the main text. INT8 post-training
quantization is reported as effectively lossless by the authors of the main
methods \citep{dettmers2022llmint8,xiao2023smoothquant} (\evC). At 4 bits, an
independent benchmark of 72 small models finds reasoning largely preserved
\citep{srivastava2025thinkslm}, and an independent benchmark of two sub-1B models
finds that methods tuned on large models transfer imperfectly, with sensitivity
depending on architecture \citep{slmquant2025} (\evB); neither establishes that
the loss grows monotonically as size falls. \gls{awq} is reported to retain more
accuracy than GPTQ at matched bit-width \citep{lin2024awq} (\evC), and mixed
schemes such as \gls{gguf} Q4\_K\_M keep the most sensitive tensors at higher
precision \citep{gerganov2023llamacpp}. In the controlled comparisons at
0.5--3.8B, post-training quantization preserved fidelity better than
post-training pruning \citep{kurz2025pruning_vs_quant,srivastava2025thinkslm}
(\evA); Section~\ref{subsec:compression_quantization} gives the scope.

Per-model, per-benchmark degradation figures are not reported consistently
enough across this cohort to tabulate: sources differ in harness, calibration
set, and group size, and few publish a matched FP16 baseline alongside the
quantized result. Measuring them under a single toolchain is one of the
verification studies proposed in
Section~\ref{subsec:challenges_verification}.

The memory consequence, by contrast, is firmly established, because it follows
from the storage format rather than from measurement.
Table~\ref{tab:memory_footprint} gives weight footprints estimated from
parameter count and effective bits per weight. The FP16 column is exactly two
bytes per parameter. The four-bit columns correspond to 4.4--5.8 effective bits
per parameter: four-bit weights plus per-group scales and zero-points, with the
embedding and output layers typically kept at higher precision, which is a
larger share of a smaller model. The reduction therefore falls short of the
nominal 4$\times$, at 2.8 to 3.6$\times$ across the table, and is smallest for
the smallest models. Qwen2.5-7B drops from 15.2~GB to 4.2--4.5~GB, and Phi-4-mini
(3.8B) from 7.6~GB to 2.2--2.4~GB, which is what moves each of them across
a hardware tier boundary. The exact figure for a given artifact is its file
size, which should be read from the artifact rather than from this table.

\begin{table}[ht]
\centering
\caption{Model-weight memory footprint by precision (GB). FP16 is two bytes per parameter; the four-bit columns are estimates at 4.4--5.8 effective bits per parameter, including the group scales and higher-precision tensors each scheme retains, which is why the reduction ($2.8$--$3.6\times$) is below the nominal 4$\times$. Activation and KV-cache memory are additional and depend on context length and batch size.}
\label{tab:memory_footprint}
\footnotesize
\begin{tabular}{lcccc}
\toprule
\textbf{Model} & \textbf{FP16} & \textbf{GPTQ-4bit} & \textbf{AWQ-4bit} & \textbf{Q4\_K\_M} \\
\midrule
Qwen2.5-7B & 15.2 & 4.2 & 4.4 & 4.5 \\
Phi-4-mini & 7.6 & 2.2 & 2.3 & 2.4 \\
Llama 3.2-3B & 6.4 & 1.9 & 2.0 & 2.1 \\
SmolLM2-1.7B & 3.4 & 1.1 & 1.1 & 1.2 \\
TinyLlama-1.1B & 2.2 & 0.7 & 0.8 & 0.8 \\
\bottomrule
\end{tabular}
\end{table}

